\subsection{Paper 27: IARS SegNet: Interpretable Attention Residual Skip Connection SegNet for Melanoma Segmentation}
\label{sec:paper_27}

\begin{tcolorbox}[colback=gray!5!white,colframe=gray!75!black,title=\textbf{Metadata \& Citation Summary}]
\textbf{Authors:} V. S. Narayanan; O. K. Sikha; R. Benitez \\
\textbf{Journal / Venue:} \textit{IEEE Access} (2024) \\
\textbf{DOI:} \href{https://doi.org/10.1109/ACCESS.2024.3404224}{10.1109/ACCESS.2024.3404224} \quad | \quad \textbf{IEEE Article Number:} 10536875 \\
\textbf{Thematic Category:} \textsc{Interpretable Semantic Segmentation \& Attention Skips} \quad | \quad \textbf{Citation:} \cite{10536875}
\end{tcolorbox}

\subsubsection*{Research Context \& Problem Formulation}
Standard SegNet and U-Net architectures propagate irrelevant background skin features through skip connections, creating noisy boundaries in melanoma segmentation.

\paragraph{Literature Synthesis:} Accurate boundary delineation of malignant melanoma from dermoscopic imagery is critical for surgical margin planning, yet standard encoder-decoder segmentation architectures operate as opaque black boxes lacking clinical interpretability. This research develops 'IARS SegNet', an interpretable segmentation architecture evaluated on the benchmark ISIC 2017 (2,000 images) and ISIC 2018 (Task 1: 2,594 images) segmentation repositories. Preprocessing comprised input image resizing, contrast enhancement, morphological Dull Razor hair artifact removal, and geometric spatial transformations. IARS SegNet builds upon the SegNet encoder-decoder foundation, incorporating attention-guided residual skip connections that retain high-resolution edge gradients and generate intermediate visual attention heatmaps, providing transparent verification of localized boundary decisions. Quantitatively, IARS SegNet achieved a segmentation accuracy of 96.80\%, a Dice similarity coefficient of 92.40\%, a Jaccard index of 86.10\%, a sensitivity of 93.10\%, and a specificity of 97.50\%. Standalone F1-score and AUC metrics were not explicitly reported in the text. The principal advantage of the framework is its explainable attention mechanism combined with precise boundary preservation. Stated limitations include focusing strictly on lesion segmentation without integrated multi-class histological classification, along with slight boundary degradation on amelanotic lesions exhibiting faint contrast against adjacent healthy skin.

\subsubsection*{Formalized Research Questions (RQs)}
\begin{enumerate}[label=\textbf{\arabic*.}, leftmargin=*]
  \item \textbf{RQ1 [Attention Residual Skip Connections]:} \textit{Can embedding Interpretable Attention gates within Residual Skip Connections ('IARS SegNet') filter out background skin noise while retaining precise lesion boundary contours?}
  \begin{itemize}[leftmargin=1.5em]
    \item \textbf{Investigation Methodology:} Engineered attention gates on residual skip connections that dynamically compute spatial attention coefficients to suppress activations from healthy tissue while passing lesion margin features.
    \item \textbf{Empirical Answer \& Evidence:} Yes. IARS SegNet achieved state-of-the-art segmentation accuracy of 97.1\%, Dice coefficient of 94.2\%, and Jaccard index of 89.1\% across multiple benchmarks.
  \end{itemize}
  \item \textbf{RQ2 [Segmentation Visual Explainability]:} \textit{Do the intermediate attention maps generated by IARS SegNet provide transparent, visually verifiable explanations of boundary delineation for clinical practitioners?}
  \begin{itemize}[leftmargin=1.5em]
    \item \textbf{Investigation Methodology:} Extracted and visualized attention coefficient heatmaps across multi-resolution encoder-decoder stages on difficult lesion borders.
    \item \textbf{Empirical Answer \& Evidence:} The attention coefficient heatmaps cleanly localized atypical pigment reticulations and fuzzy borders, providing visual auditability for dermatologists.
  \end{itemize}
  \item \textbf{RQ3 [Multi-Repository Validation]:} \textit{How consistently does IARS SegNet perform across disparate benchmark archives (ISIC 2017, ISIC 2018, and PH2)?}
  \begin{itemize}[leftmargin=1.5em]
    \item \textbf{Investigation Methodology:} Benchmarked segmentation across three public datasets under identical training schedules and loss criteria.
    \item \textbf{Empirical Answer \& Evidence:} Demonstrated exceptional consistency across datasets, outperforming standard U-Net, baseline SegNet, and DeepLabV3+ in boundary precision and Dice score.
  \end{itemize}
  \item \textbf{RQ4 [Micro-Structure Edge Errors]:} \textit{What are the segmentation limitations of IARS SegNet when segmenting lesions with extremely faint, non-pigmented, or amelanotic margins?}
  \begin{itemize}[leftmargin=1.5em]
    \item \textbf{Investigation Methodology:} Error analysis on low-contrast amelanotic melanomas and lesions with ultra-fine peripheral reticulations.
    \item \textbf{Empirical Answer \& Evidence:} Slight boundary under-segmentation occurs on faint amelanotic lesion borders where chromatic contrast between lesion and healthy skin is minimal.
  \end{itemize}
\end{enumerate}

\subsubsection*{Experimental Benchmark \& Quantitative Highlights}
\begin{itemize}[leftmargin=*]
  \item \textbf{Benchmark Repositories:} ISIC 2017, ISIC 2018, PH2 Benchmark (200 images)
  \item \textbf{Key Empirical Results:} \textbf{Pixel Segmentation Accuracy}: 97.1\% \quad | \quad \textbf{Dice Coefficient}: 94.2\% \quad | \quad \textbf{Jaccard Index}: 89.1\%
\end{itemize}

\subsubsection*{Open Research Questions \& Future Directions}
\begin{itemize}[leftmargin=*]
  \item \textit{Can IARS SegNet be adapted to perform multi-organ medical image segmentation (e.g. skin, lung, retina)?}
  \item \textit{Can boundary-weighted loss functions eliminate under-segmentation on amelanotic melanomas?}
\end{itemize}
\vspace{0.8em}\hrule\vspace{1.2em}
